\noindent
Genes associated with educational attainment causally improve labour market income, but the economic mechanism behind this relationship is not clear.
Using quasi-random variation in genetic inheritance across siblings in the UK Biobank, I estimate the causal effect of the Education PolyGenic Index (Ed PGI) on education and income, controlling for imputed parental PGI values to isolate the random inheritance component.
A one standard deviation increase in the Ed PGI raises completed education by 0.5 years and later life income by around 5 percent (replicating the main estimates in \citealt{carvalho2024genetics}).
I then decompose this total genetic income effect into an indirect channel operating through education years and a residual direct effect, using a causal mediation framework.
Unlike existing model-based decompositions, the approach is design-based: the one genuine dimension of uncertainty --- the causal return to a year of education --- is handled transparently through a sensitivity analysis benchmarked against quasi-experimental estimates from the economics literature for Britain.
At correlational returns of around 6 percent, 65 to 75 percent of the total genetic income effect operates through the education channel; at plausible causal returns of around 10 percent, the direct genetic effect closes towards zero.
%Education-linked genes earn their labour market return primarily by inducing more schooling, not through direct biological pathways that bypass education.
The institutions governing access to education are therefore the primary mechanism through which genetic endowment shapes income inequality, and the primary policy lever through which that inequality could be addressed.



\vspace{0.25cm}
\noindent
\textbf{Keywords:}
Education,
Mediation,
Polygenic indicies,
Nature/nurture.

\vspace{0.1cm}
\noindent
\textbf{JEL Codes:}
D31, D91, I24, J24, Z00.
